\clearpage

\section{总结与展望}

\subsection{论文总结}

本文围绕大规模接入物联网中的状态更新问题，研究了面向 AoI 的随机接入机制设计与性能分析。针对现有工作对持续碰撞重视不足、对 AoI 违规概率讨论较少的情况，本文提出了一种基于等待的随机接入退避机制，并将其分别放到时隙 ALOHA 和标准 CSMA 两类经典框架下进行分析，得到了平均 AoI 与 AoI 违规概率的闭式结果，并通过仿真验证了相关结论。全文的主要工作和结论可以归纳为以下四点：
\begin{enumerate}[wide]
    \item \textbf{提出一种适用于大规模接入场景的等待机制。}考虑到高负载条件下碰撞会持续累积，并进一步拉长有效更新之间的间隔，本文没有只把注意力放在单次碰撞上，而是把节点在成功更新与碰撞之后的再次接入过程一并纳入设计。通过让节点在成功更新后适度等待、在碰撞后执行随机退避，系统中的竞争强度能够得到调节，有效更新在时间轴上的分布也更均匀。

    \item \textbf{围绕时隙 ALOHA 建立等待机制的理论分析框架。}针对大规模网络中节点状态强耦合、系统难以直接解析的问题，本文在稳态解耦的近似下构建了单节点马尔可夫链模型，求解了相应的稳态方程，并据此推导出平均 AoI 与 AoI 违规概率的闭式表达式，进一步分析了等待窗口、网络规模等参数对系统时效性的影响。

    \item \textbf{围绕标准 CSMA 建立等待机制的理论分析框架。}针对大规模网络中节点状态强耦合，以及载波侦听、退避冻结导致的变长系统时隙所带来的分析困难，本文同样在稳态解耦的近似下构建了单节点马尔可夫链模型，求解了相应的稳态方程，并将稳态结果与变长时隙的统计特性结合起来，进一步推导出平均 AoI 与 AoI 违规概率的闭式表达式，分析了等待窗口、网络规模等参数对系统时效性的影响。

    \item \textbf{通过理论与仿真对比揭示等待机制的作用规律。}理论结果与仿真结果基本一致，说明本文建立的模型能够较好地刻画系统性能。进一步的对比表明，在中高负载尤其是海量设备接入的场景下，适度等待不会必然带来信息时效性的恶化，反而有助于减轻信道竞争、降低碰撞概率，并同时改善平均 AoI 与 AoI 违规概率。
\end{enumerate}

综合全文的分析与仿真结果可以看出，等待机制的关键在于调节节点再次接入信道的时机，从而缓和大规模接入场景下的过度竞争。在物联网状态更新业务中，这种调节有助于使有效更新在时间上分布得更均匀，进而改善系统的信息时效性。本文得到的相关结论，也可为后续面向 AoI 的随机接入协议设计与优化提供参考。

\subsection{研究展望}

本文的工作还建立在若干理想化假设之上，因此仍有一些问题值得继续往前推进。结合本文的分析边界，后续研究可以从以下几个方向展开：
\begin{enumerate}[wide]
    \item \textbf{进一步纳入更真实的到达与排队因素。}本文主要采用按需更新的建模方式，重点分析接入竞争过程对 AoI 的影响，还没有把随机业务到达、缓存队列和服务过程统一放进同一个模型。后续如果能在排队论框架下把这些因素结合起来，就能更细致地考察等待机制在随机到达流量、有限缓存以及多业务混合场景中的表现。

    \item \textbf{推广到更复杂的信道和网络条件。}本文关注的是单跳共享信道环境，对信道衰落、隐藏终端、节点异构和多跳拓扑等因素涉及不多。把等待机制放到这些更复杂的无线环境中，一方面可以检验现有结论是否稳健，另一方面也有助于看清不同网络条件下信息时效性的变化规律。

    \item \textbf{继续讨论参数选择与自适应设计。}虽然本文已经给出了平均 AoI 与 AoI 违规概率的解析结果，但等待窗口、退避参数和系统负载之间究竟如何匹配，仍有进一步挖掘的空间。后续可以围绕这一问题建立参数优化模型，并探索根据网络负载或碰撞状态动态调整等待参数的自适应机制。

    \item \textbf{结合实际协议平台开展验证。}本文目前的结论主要来自理论分析和数值仿真。若能进一步借助基于标准 CSMA 的实验平台、网络仿真器或软硬件原型系统实现等待机制，就能更具体地评估它的工程开销、部署难点和真实环境下的性能收益。
\end{enumerate}
